\subsection{Редукция задачи Коши для УКБС}
\[
	u(x, t) = U(x, t) + V(x, t) + W(x, t),
\] где $U(x, t), V(x, t), W(x, t)$ - решения следующих задач Коши для УКБС:

\begin{align*}
	                  & U_{tt} = a^2 U_{xx}, \quad x \in \mathbf{R}, \ t>0,          \\
	\text{НУ: } \quad & U|_{t=0}=u_0(x),\ U_t|_{t=0}=0,\ x \in \mathbf{R}            \\
	                  & V_{tt} = a^2 V_{xx}, \quad x \in \mathbf{R}, \ t>0,          \\
	\text{НУ: } \quad & V|_{t=0}=0,\ V_t|_{t=0}=u_1(x),\ x \in \mathbf{R}            \\
	                  & W_{tt} = a^2 W_{xx} + f(x,t), \quad x \in \mathbf{R}, \ t>0, \\
	\text{НУ: } \quad & W|_{t=0}=W_t|_{t=0}=0,\ x \in \mathbf{R}
\end{align*}
В формуле Даламбера:

\begin{itemize}[nosep, leftmargin=1.5cm]
	\item $
		      U(x, t) = \frac{1}{2}[u_0(x-at) + u_0(x+at)]
	      $

	\item $
		      V(x, t) = \frac{1}{2a} \int_{x-at}^{x+at}u_1(\xi)\,dx
	      $

	\item $
		      W(x, t) = \frac{1}{2a} \int_0^t \,dt \int_{x-a(t-\tau)}^{x+a(t-\tau)} f(\xi, \tau) \,d\tau
	      $
\end{itemize}

\subsection{Принцип Дюамеля}

\begin{lemma}
	Пусть функция u(x, t, $\tau$) при всех $\tau \ge 0$ является решением ЗК для ОУКБС:
	\[
		u_{tt}=a^2u_{xx},\ x \in \mathbf{R},\ t > \tau \ \\
		\text{НУ: }u|_{t=\tau}=0,\ u_t|_{t=\tau}=f(x, \tau),\ x \in \mathbf{R}
	\]
	Тогда функция $v(x, t) = \int_0^t u(x, t, \tau)d\tau$ является решением ЗК для УКБС:
	\[
		v_{tt}=a^2v_{xx} + f(x, t),\ x \in \mathbf{R},\ t > 0 \ \\
		\text{НУ: }v|_{t=0}=u_t|_{t=0}=0,\ x \in \mathbf{R}
	\]
\end{lemma}

\begin{proof}
	В самом деле, если $v(x, t) = \int_0^tu(x, t, \tau)d\tau$, где $u(x, t, \tau)$ при $\forall \tau \ge 0$ является решением ЗК для ОУКБС $(1.0),\ (2.f)$, то
	\[
		v_{t=0} = v(x, 0) = \int_0^0u(x, t, \tau) = 0,\ x \in \mathbf{R}
	\]

	\[
		v_t = u(x, t, t) + \int_0^tu_t(x, t, \tau) d\tau \Rightarrow v_t|_{t=0}=0+0=0
	\]
	Проверили, что решение $v(x, t)$ удовлетворяет НУ.

	Далее
	\[
		v_{tt} = u_t(x, t, t) + \int_0^t u_{tt}(x, t, \tau) d\tau = f(x, t) + a^2 \int_0^t u_{xx} d\tau = f(x, t) + a^2 \frac{\partial^2}{\partial x^2} \int_0^tu(x, t, \tau) d\tau = a^2v_{xx} + f(x, t)
	\]
\end{proof}

Таким образом, в соответствии с принципом Дюамеля, решение $W(x, t)$ ЗК $(1.f),\ (2.0)$ для УКБС имеет вид:
\[
	W(x, t) = \int_0^tu(x, t, \tau) d\tau,
\]где $u(x, t,\tau)$ --- решение ЗК $\forall \tau \ge 0 (1.0), (2.f)$ для ОУКБС. А по (ФДО):
\[
	u(x, t, \tau) = \frac{1}{2a} \int_{x-a(t-\tau)}^{x+a(t-\tau)} f(\xi, \tau) d\xi \Rightarrow
\]
\[
	W(x, t) = \frac{1}{2a}\int_0^T d\tau \int_{x-a(t-\tau)}^{x+a(t-\tau)} f(\xi, \tau) d\xi
\]

\subsection{Интеграл энергии}

Рассмоитри простейшую СЗ для ОУКС длины l, закрепленной на концах
\begin{align*}
	                       & u_{tt} = a^2 u_{xx}, \quad x \in (0,l),\ t > 0 \tag{1.0}      \\
	\textbf{НУ: } \quad    & u|_{t=0} = u_0(x),\ u_t|_{t=0} = u_1(x),\ x \in [0,l] \tag{2} \\
	\textbf{ОГУ1Т: } \quad & u|_{x=0} = u|_{x=l} = 0,\ t \ge 0 \tag{3.1.0}                 \\
	\textbf{ОУС: } \quad   & u_0(0) = u_0(l) = u_1(0) = u_1(l) = 0
\end{align*}
\[
	u(x, t) \in C^2(G) \cap C^1(\overline{G}),
\]
\[
	u_0(x) \in C^2([0, l]),\ u_1(x) \in C^1([0, l]),
\]
где $G=(0, l) \times (0, +\infty)$ и $\overline{G}=[0, l] \times [0, +\infty)$.

И рассмотрим интеграл
\begin{align*}
	0 & = \int_0^l \frac{\partial u}{\partial t} Lu\ dx = \int_0^l\left[\rho \frac{\partial u}{\partial t} \frac{\partial^2 u}{\partial t^2} -T \frac{\partial u}{\partial t} \frac{\partial^2 u}{\partial x^2} \right]dx = \int_0^l \frac{\rho}{2} \frac{\partial}{\partial t}\left(\frac{\partial^2 u}{\partial t^2} \right)dx - \int_0^l T \frac{\partial u}{\partial t} d\left(\frac{\partial u}{\partial x} \right)                                                                      \\
	  & = \int_0^l \rho \frac{\partial u}{\partial t} \frac{\partial^2 u}{\partial t^2} - T \frac{\partial u}{\partial t} \left. \frac{\partial u}{\partial x}\right|_0^l + \int_0^l \frac{\partial u}{\partial x} \left(d \left(T \frac{\partial u}{\partial t} \right) \right) = \int_0^l \frac{\rho}{2} \frac{\partial}{\partial t} \left( \frac{\partial u}{\partial t} \right)^2 dx + \int_0^l \frac{T}{2} \frac{\partial}{\partial t} \left( \frac{\partial u}{\partial t} \right)^2 dx \\
	  & = \left(\text{т.к.} \frac{\partial u}{\partial x} \frac{\partial^2 u}{\partial x \partial t} =
	\frac{1}{2} \frac{\partial }{\partial t} \left( \frac{\partial u}{\partial x} \right)^2
	\right) = \int_0^l \frac{\partial}{\partial t} \left[ \frac{\rho}{2} \left( \frac{\partial u}{\partial t} \right)^2 + \frac{T}{2}\left( \frac{\partial u}{\partial x} \right)^2 \right] dx                                                                                                                                                                                                                                                                                                \\
	  & = \frac{d}{dt} \int_0^l \left[ \frac{\rho}{2} \left( \frac{\partial u}{\partial t} \right)^2 + \frac{T}{2}\left( \frac{\partial u}{\partial x} \right)^2 \right]dx = 0 \Rightarrow
\end{align*}

\[
	E(t) = \int_0^l \left[ \frac{\rho}{2} \left( \frac{\partial u}{\partial t} \right)^2 + \frac{T}{2}\left( \frac{\partial u}{\partial x} \right)^2 \right]dx = E(0) = const \text{ --- интеграл энергии колебания струны}
\]
При этом
\[
	\omega(x, t) = \frac{\rho}{2} \left( \frac{\partial u}{\partial t} \right)^2 + \frac{T}{2}\left( \frac{\partial u}{\partial x} \right)^2 \text{ --- плотность энергии колебания струны}
\]

\subsection{Смешанная задача для УКС с ГУ1Т(длина l)}

\begin{align*}
	                      & u_{tt} = a^2 u_{xx} + f(x, t), \quad x \in (0,l),\ t > 0 \tag{1.f} \\
	\textbf{НУ: } \quad   & u|_{t=0}=u_0(x),\ u_t|_{t=0}=u_1(x),\ x \in [0,l] \tag{2}          \\
	\textbf{ГУ1Т: } \quad & u|_{x=0}=\mu(t),\ u|_{x=l}=\nu(t),\ t \ge 0 \tag{3.1}              \\
	\textbf{УС: } \quad   & u_0(0)=\mu(0),\ u_0(l)=\nu(0), \tag{4}
\end{align*}

Теперь докажем единственность решения смешанной задачи(СЗ) для УКС(уравнения колебания струны) с ГУ1Т(граничными условиями 1 типа) $(1.f),\ (2),\ (3.1),\ (4)$.

\begin{theorem}[единственность решения]
	Если решение СЗ для УКС длины l с ГУ1Т $(1.f),\ (2),\ (3.1),\ (4)$ существует в рассматриваемых классах функций, то оно единственно
\end{theorem}

% \begin{proof}
% Пусть $\tilde u(x, t)$ и $\hat u$ --- решения поставленной задачи, тогда 
% \[
% E(t) = \int_0^l \left[ \frac{\rho}{2} \left( \frac{\partial u}{\partial t} \right)^2 + \frac{T}{2}\left( \frac{\partial u}{\partial x} \right)^2 \right]dx = E(0) = const
% \]
% и функция $u(x, t) = \tilde{u}(x, t) - \hat{u}(x, t)$ --- является решением СЗ для однородного УКС длины l, закрепленной на концах, удовлетворяющей однородным НУ:
% \[
% Lu = \rho u_{tt} - T u_{xx} = 0,\ x \in (0, l), t > 0\quad (1.0) \ \\
% \textbf{НУ: } u|_{t=0} = u_t|_{t=0} = 0 ,\ x \in [0, l]\quad (2.0) \ \\
% \textbf{ОГУ1Т: } u|_{x=0} = u_t|_{x=l} = 0 ,\ t \ge 0 \quad (3.1.0)
% \]
% \end{proof}